\scanpage{63}
\mspage{62}
% E: text|upper left|et tout chemin
et tout chemin \inserted{$c$} définit des iso. d'ens à deux éléments
% E: diagram|upper left|omega epsilon_s epsilon_a
\diagram{\begin{tikzcd}[column sep=large,row sep=small]
\underline\omega(\unsure{F|D})\simeq\underline\omega(D) \arrow[r,"\underline\omega_X(c)"] & \underline\omega(\unsure{F'|D'})\simeq\underline\omega(D')\\
\varepsilon_s(F)=\varepsilon_s(D) \arrow[r] & \varepsilon_s(F')\simeq\varepsilon_s(D')\\
\varepsilon_a(F)\unsure{\simeq|=}\varepsilon_a(D) \arrow[r] & \varepsilon_a(F')\simeq\varepsilon_a(D')
\end{tikzcd}}
% E: text|left upper middle|Ces ens. ne sont pas indépendants
Ces ens. ne sont pas indépendants, chaque paire des deux détermine le troisième, par la relation
% E: display|left upper middle|omega(F) epsilon_s(F) wedge epsilon_a(F)
\[\underline\omega(F)\simeq\varepsilon_s(F)\wedge\varepsilon_a(F)\]
% E: text|left middle|et les deux relations
et les deux relations \unsure{circulaires|analogues} équivalentes.

Au niveau des caractères correspondants sur $\pi_1(X,F)$, cela donne
% E: diagram|left middle|chi_X(F)=chi_s(F)chi_a(F)
\diagram{\begin{tikzcd}[row sep=small]
\boxed{\chi_X(F)=\chi_s(F)\chi_a(F)} \arrow[d,phantom,"\rotatebox{90}{$=$}" description]\\
\chi_X(D)=\chi_s(D)\chi_a(D) \arrow[d,phantom,"\rotatebox{90}{$=$}" description]\\
\underline\omega(D)
\end{tikzcd}}
% E: text|left middle|ou de façon plus symétrique
ou de façon plus symétrique
% E: display|left middle|chi_X(F)chi_s(F)chi_a(F)=1
\[\chi_X(F)\chi_s(F)\chi_a(F)=1.\]
% E: text|lower left|En plus de ces trois caractères
En plus de ces trois caractères, il y en a deux autres, $\chi_{\mathcal X}={}$caractère d'\unsure{orientabilité|orientation} de $\mathcal X$ (ne dépendant pas de la structure cellulaire), et $\chi_c$ (``combinatoire''), qui est donné par la parité \inserted{\unsure{du tout|de tout}}. Le système local \deleted{\unsure{associé est les|associé les}} \inserted{correspondant est défini par les} $\underline\omega(F)$, mais avec des identifications de transition \uline{opposées} à celles données par \deleted{\unsure{identifications}} le transport parallèle sur $X$
% E: margin note|diagonal centre with arrow to parity phrase|Attention dans le
\margin{Attention, dans le \unsure{pointé|point}, \uline{est faux}, il faut compter pour \uline{deux} le franchissement d'un côté \unsure{except\textsuperscript{l}|except.} (qui \unsure{accroche|accrocherait}), \unsure{soit décomposer en|décomposé en} \uline{deux} opérations.}
% E: text|upper right|défini par les omega_X(c)
défini par les $\underline\omega_X(c)$ --- donc on aura
% E: display|upper right|omega_Xcal(c)=p(c)omega_X(c)
\[\underline\omega_{\mathcal X}(c)=\unsure{\mathrm{par}|p}(c)\omega_X(c)\]
% E: text|upper right|où p(c) désigne la parité
\unsure{où} $\unsure{\mathrm{par}|p}(c)$ désigne la parité de $c$.
\deleted{En prenant} Le système local associé à $\chi_c$ a \deleted{formé de} comme fibres les ens. triviaux $\{\pm1\}$, mais les morphismes de transition donnés par les $\unsure{\mathrm{par}|p}(c)$.

On en a donc
% E: display|right middle|chi_Xcal=chi_o chi_c
\[\boxed{\chi_{\mathcal X}=\chi_{\unsure{o|X}}\chi_c}\]
% E: margin note|right edge alongside boxed formula|On a chi_c=1
\margin{On a $\chi_c=1$ \unsure{cf p\textsuperscript{o}|si $\Gamma$} \unsure{suivante!}}
% E: text|right middle|Donc compte tenu de ces relations
Donc, compte tenu de ces relations, on dispose de trois caractères fondamentaux \inserted{\unsure{privilégiés}}, que je choisis
% E: display|right lower middle|chi_X chi_s chi_c
\[\chi_X,\chi_s,\chi_c\]
% E: text|right lower middle|donc les deux autres sont
\unsure{donc} les deux autres sont
% E: display|right lower middle|chi_a=chi_X chi_s
\[\chi_a=\chi_X\chi_s\quad\text{et surtout}\quad\chi_{\mathcal X}=\chi_X\chi_c\]
% E: text|bottom right|Il y a lieu d'expliciter
Il y a lieu d'expliciter les valeurs des caractères précédents, si faire se peut, ainsi que celles des représentations
% E: display|bottom right|u pi_1(X F_D) to Aut(Pol(D))
\[u:\pi_1(X,F_D)\longrightarrow\operatorname{Aut}(\operatorname{Pol}(D)),\]
